% Unit 27 — Reverse quotient rule
\input{technique-note.tex}
\title{Integration Bee --- Unit 27:\\ Reverse Quotient Rule}
\author{UK Uni Integration Bee prep\\ The quotient rule run backwards}
\date{}
\begin{document}
\maketitle
\section*{Status}
Tags: \texttt{yang:4.6}; techniques: \texttt{reverse-quotient}.
The quotient rule says $(f/g)'=(f'g-fg')/g^2$. Integrating the right-hand side returns $f/g$. Yang \S4.6 collects contest shapes of that right-hand side. Unit 07 is the product-rule twin. You do not need a new calculus fact.
\section{Official anchors}
\begin{itemize}
\item The quotient rule, any calculus text. Unit 07 for the product-rule version.
\item Syllabus \S1.
\item Yang Topics Guide 2nd Ed \S4.6 Reverse Quotient Rule.
\end{itemize}
\section{Learning goals}
\begin{enumerate}
  \item Read a numerator $f'g-fg'$ and a denominator $g^2$ as $(f/g)'$.
  \item Guess $f/g$ when the denominator is a square, then differentiate to check.
\end{enumerate}

\section{Worked --- \texttt{U27-D01}}
\[
\int\frac{x\cos x-\sin x}{x^2}\,dx.
\]
Here $f=\sin x$ and $g=x$. Then $f'g-fg'=x\cos x-\sin x$ and $g^2=x^2$. The antiderivative is $\sin x/x$. That is the stored answer.

\texttt{U27-D04} is $( \sin x-x\cos x)/\sin^2 x$, so $f=x$ and $g=\sin x$, and the answer is $x/\sin x$. \texttt{U27-D07} is $(x\sec^2 x-\tan x)/x^2$, so $f=\tan x$ and $g=x$. \texttt{U27-D08} is $(x^2\cos x-2x\sin x)/x^4=(f'g-fg')/g^2$ with $f=\sin x$ and $g=x^2$, answer $\sin x/x^2$.

\section{Exponentials and logarithms}
\texttt{U27-D02} is $e^x(x-1)/x^2$. Take $f=e^x$ and $g=x$:
\[
f'g-fg'=e^x\cdot x-e^x\cdot 1=e^x(x-1).
\]
The antiderivative is $e^x/x$. \texttt{U27-D06} is the same with $g=x+1$, answer $e^x/(x+1)$.

\texttt{U27-D03} is $(1-\ln x)/x^2$. Take $f=\ln x$ and $g=x$. Then $f'g-fg'=1-\ln x$, so the answer is $\ln x/x$. \texttt{U27-D05} flips the sign and stores $-\ln x/x$.

\section{Faded check}
Differentiate the proposed answer before you box it. If the denominator of the derivative is not the square in the integrand, the guess had the wrong $g$. A numerator that is $f'$ with no minus term is a logarithmic derivative, unit 24, not this one.

\section{Drill pointers}
See \texttt{drills.tex} / \texttt{drills.md}. Feed: \texttt{verify/drills\_unit27.json}.

\section*{On a miss}
Teach this file. Name $f$ and $g$, then differentiate $f/g$. Do not invent an answer that is not in \texttt{drills.md}.
\end{document}
